\documentclass[11pt,a4paper,oneside]{book}
\usepackage[utf8]{inputenc}
\usepackage[italian]{babel}
\usepackage{hyperref}
\usepackage{listings}
\usepackage{xcolor}
\usepackage{graphicx}
\usepackage{geometry}
\usepackage{titlesec}
\usepackage{fancyhdr}
\usepackage{tcolorbox}

\geometry{margin=2.5cm, headheight=15pt}

% Stili per il codice
\definecolor{codegreen}{rgb}{0,0.6,0}
\definecolor{codegray}{rgb}{0.5,0.5,0.5}
\definecolor{codepurple}{rgb}{0.58,0,0.82}
\definecolor{backcolour}{rgb}{0.96,0.96,0.96}

\lstdefinestyle{mystyle}{
    backgroundcolor=\color{backcolour},   
    commentstyle=\color{codegreen},
    keywordstyle=\color{magenta},
    numberstyle=\tiny\color{codegray},
    stringstyle=\color{codepurple},
    basicstyle=\ttfamily\footnotesize,
    breakatwhitespace=false,         
    breaklines=true,                 
    captionpos=b,                    
    keepspaces=true,                 
    numbers=left,                    
    numbersep=5pt,                  
    showspaces=false,                
    showstringspaces=false,
    showtabs=false,                  
    tabsize=2
}
\lstset{style=mystyle}

% Intestazioni e piè di pagina
\pagestyle{fancy}
\fancyhf{}
\fancyhead[L]{Documentazione Tecnica Gulliver 2026}
\fancyhead[R]{\thepage}
\fancyfoot[C]{Confidenziale - Solo per uso interno}

\title{
    \vspace{2cm}
    \textbf{\Huge Documentazione Tecnica Estesa}\\
    \vspace{0.5cm}
    \Large Sistema Web Integrato: Gulliver Web 2026\\
    \vspace{2cm}
    \includegraphics[width=0.3\textwidth]{logo-example} % Se vuoi puoi aggiungere il logo
}
\author{\textbf{Sviluppato da Antigravity AI}}
\date{Maggio 2026}

\begin{document}

\maketitle
\newpage
\tableofcontents
\newpage

\chapter{Introduzione al Progetto}
\section{Obiettivi}
Il progetto "Gulliver Web 2026" nasce dall'esigenza di modernizzare l'infrastruttura digitale della lista studentesca Gulliver UNIVPM. Gli obiettivi principali sono:
\begin{itemize}
    \item Centralizzare la gestione dei form esterni (Tally, Google Forms).
    \item Fornire un accesso rapido e strutturato all'archivio appunti.
    \item Garantire un'esperienza utente premium e responsive su ogni dispositivo.
    \item Semplificare il passaggio di consegne tecnico tra le generazioni di studenti.
\end{itemize}

\section{Filosofia Tecnica}
Il sito adotta una filosofia \textbf{Serverless First}. Non esiste un server fisico o una VPS da gestire; tutto gira su infrastruttura gestita (Vercel) e database edge (Upstash). Questo annulla i costi di manutenzione sistemistica e garantisce uptime elevatissimo.

\chapter{Architettura Software (Next.js 15)}
\section{Il Framework}
Next.js 15 è il framework React più avanzato al mondo. Utilizza il paradigma dell'\textbf{App Router}.
\subsection{Server Components vs Client Components}
La maggior parte delle pagine sono \textbf{Server Components}, il che significa che vengono renderizzate sul server e inviate all'utente già pronte. I \textbf{Client Components} (indicati con \texttt{'use client'}) vengono usati solo dove serve interattività (es. pulsanti, filtri di ricerca, notifiche).

\section{Struttura delle Cartelle}
\begin{itemize}
    \item \texttt{/src/app}: Contiene tutte le rotte (pagine) del sito.
    \item \texttt{/src/app/admin}: Logica del pannello di controllo.
    \item \texttt{/src/app/api}: Endpoint serverless per Redis e Google Sheets.
    \item \texttt{/src/components}: Elementi riutilizzabili (Header, Footer, Cards).
    \item \texttt{/public}: Asset statici (immagini, loghi, PDF).
\end{itemize}

\chapter{Infrastruttura di Rete e DNS}
\section{Configurazione Cloudflare}
Cloudflare agisce come scudo (WAF) e come gestore DNS.
\subsection{Configurazione SSL/TLS (Critico)}
È \textbf{fondamentale} che la modalità SSL sia impostata su \textbf{Full} o \textbf{Full (Strict)}. 
\textbf{Perché?} Vercel richiede HTTPS. Se Cloudflare è su "Flexible", proverà a connettersi a Vercel in HTTP; Vercel risponderà con un redirect a HTTPS, creando un ciclo infinito (Infinite Redirect Loop) che rende il sito inaccessibile.

\section{Gestione dei Sottodomini}
\begin{itemize}
    \item \textbf{forms.gulliverancona.it}: Gestisce i reindirizzamenti.
    \item \textbf{admin.gulliverancona.it}: Porta al pannello di controllo.
    \item \textbf{www.gulliverancona.it}: Il sito principale.
\end{itemize}

\chapter{Gestione dei Form e Redirect}
\section{Il Flusso di Redirect}
Quando un utente accede a \texttt{forms.gulliverancona.it/slug}, la rotta dinamica \texttt{/f/[slug]} intercetta la richiesta.
\begin{enumerate}
    \item Viene estratto lo \texttt{slug} dall'URL.
    \item Si effettua una chiamata \texttt{GET} a Redis per la chiave \texttt{form:slug}.
    \item Se trovata, si verifica lo stato:
    \begin{itemize}
        \item \texttt{active}: Redirect 302 a Tally.
        \item \texttt{suspended}: Mostra pagina di errore custom.
    \end{itemize}
\end{enumerate}

\section{Database Redis (Upstash)}
Usiamo Redis perché è istantaneo (chiave-valore). La struttura dati è semplice JSON salvato sotto chiavi univoche. Questo permette di gestire migliaia di redirect con latenza vicina allo zero.

\chapter{Sistema Archivio Appunti}
\section{Integrazione Google Sheets}
Il sistema non salva gli appunti nel proprio database, ma li legge "al volo" da un Google Sheet. Questo permette ai responsabili appunti di continuare a usare Excel/Sheets senza dover imparare un nuovo software.

\section{Logica di Parsing CSV}
L'endpoint \texttt{/api/appunti} scarica il foglio in formato CSV. Poiché il CSV può contenere virgole all'interno dei campi (es. nomi professori o descrizioni), abbiamo implementato un parser custom (\texttt{parseCsvLine}) che gestisce correttamente le virgolette.

\chapter{Sicurezza e Autenticazione}
\section{Portale Admin}
L'accesso all'admin è protetto da una password salvata in \texttt{ADMIN\_PASSWORD}. 
\textbf{Nota di sicurezza}: Non vengono salvati utenti singoli. La password è globale per la lista. Al login viene creato un cookie di sessione criptato.

\section{Variabili d'Ambiente (.env)}
Le variabili sensibili non sono nel codice. Devono essere impostate nel pannello Vercel (Settings -> Environment Variables).
\begin{itemize}
    \item \texttt{UPSTASH\_REDIS\_REST\_URL}
    \item \texttt{UPSTASH\_REDIS_REST\_TOKEN}
    \item \texttt{ADMIN\_PASSWORD}
    \item \texttt{NEXT\_PUBLIC\_APPUNTI\_SHEET\_ID}
\end{itemize}

\chapter{Design e UI/UX}
\section{Il Sistema di Design}
Il sito usa un tema \textbf{Dark Premium} basato su:
\begin{itemize}
    \item \textbf{Glassmorphism}: Effetto vetro sfocato (\texttt{backdrop-filter: blur}).
    \item \textbf{Aesthetics}: Colori vibranti ma desaturati per non affaticare la vista.
    \item \textbf{Micro-animations}: Transizioni fluide su ogni hover e click.
\end{itemize}

\chapter{Guida alla Manutenzione}
\section{Aggiornare il Sito}
Ogni volta che fai un \texttt{git push} sul branch \texttt{main}, Vercel avvia automaticamente un nuovo deploy. Se la build ha successo, il sito si aggiorna in 1-2 minuti.

\section{Cosa fare se il sito non carica?}
\begin{enumerate}
    \item Controlla i log su Vercel.
    \item Verifica che le credenziali Redis su Upstash non siano scadute.
    \item Assicurati che il Google Sheet degli appunti sia ancora "Pubblico su Web" (chiunque abbia il link può visualizzare).
\end{enumerate}

\appendix
\chapter{Codice Sorgente Rilevante}
In questa sezione riportiamo frammenti di codice critici per la comprensione del sistema.
\section{Parser CSV Custom}
\begin{lstlisting}[language=javascript]
function parseCsvLine(line) {
  const result = [];
  let current = '';
  let inQuotes = false;
  for (let i = 0; i < line.length; i++) {
    const ch = line[i];
    if (ch === '"') inQuotes = !inQuotes;
    else if (ch === ',' && !inQuotes) {
      result.push(current.trim());
      current = '';
    } else current += ch;
  }
  result.push(current.trim());
  return result;
}
\end{lstlisting}

\chapter{Contatti e Passaggio di Consegne}
Il responsabile tecnico attuale è Lorenzo Califano (\texttt{cali@gulliverancona.it}). In caso di successione, assicurarsi che il nuovo responsabile abbia accesso a:
\begin{itemize}
    \item Account GitHub (Organization: rappresentantigullivering).
    \item Pannello Vercel/Netlify.
    \item Console Upstash.
    \item Pannello DNS Cloudflare.
\end{itemize}

\end{document}
